\section{Equality}
\label{sec:equality}

% Derive equality and its properties - I never liked the idea of just randomly adding equality as soon as you build out ZFC.

\subsection{Definition}
\[a = b \coloneqq \forall c: ((c \in a \Leftrightarrow c \in b) \wedge (a \in c \Leftrightarrow b \in c))\]

\subsection{Properties}
\begin{align*}
    \byrule{d \in a \Leftrightarrow d \in a}{{\Leftrightarrow}R}\\
    \byrule{a \in d \Leftrightarrow a \in d}{{\Leftrightarrow}R}\\
    \byrule{(d \in a \Leftrightarrow d \in a)
        \wedge (a \in d \Leftrightarrow a \in d)}{{\wedge}I}\\
    \byrule{\forall d: ((d \in a \Leftrightarrow d \in a)
        \wedge (a \in d \Leftrightarrow a \in d))}{{\forall}I}\\
    \bydefinition{a = a}\\
    \byrule{\forall a: a = a}{{\forall}I}
\end{align*}
\[\vdash \forall a: a = a \quad (\text{reflexivity})\]

\begin{align*}
    \assumption[1]{a = b}\\
    \bydefinition[1]{\forall c: ((c \in a \Leftrightarrow c \in b) \wedge (a \in c \Leftrightarrow b \in c))}\\
    \byrule[1]{(d \in a \Leftrightarrow d \in b) \wedge (a \in d \Leftrightarrow b \in d)}{{\forall}E}\\
    \byrule[1]{d \in a \Leftrightarrow d \in b}{{\wedge}E}\\
    \byrule[1]{d \in b \Leftrightarrow d \in a}{{\Leftrightarrow}S}\\
    \byrule[1]{a \in d \Leftrightarrow b \in d}{{\wedge}E}\\
    \byrule[1]{b \in d \Leftrightarrow a \in d}{{\Leftrightarrow}S}\\
    \byrule[1]{(d \in b \Leftrightarrow d \in a) \wedge (b \in d \Leftrightarrow a \in d)}{{\wedge}I}\\
    \byrule[1]{\forall d: ((d \in b \Leftrightarrow d \in a) \wedge (b \in d \Leftrightarrow a \in d))}{{\forall}I}\\
    \bydefinition[1]{b = a}\\
    \byrule{a = b \Rightarrow b = a}{{\Rightarrow}I}\\
    \byrule{\forall a, b: (a = b \Rightarrow b = a)}{{\forall}I}\\
\end{align*}
\[\vdash \forall a, b: (a = b \Rightarrow b = a) \quad (\text{symmetry})\]

\begin{align*}
    \assumption[1]{a = b \wedge b = c}\\
    \byrule[1]{a = b}{{\wedge}E}\\
    \bydefinition[1]{\forall d: ((d \in a \Leftrightarrow d \in b) \wedge (a \in d \Leftrightarrow b \in d))}\\
    \byrule[1]{(f \in a \Leftrightarrow f \in b) \wedge (a \in f \Leftrightarrow b \in f)}{{\forall}E}\\
    \byrule[1]{f \in a \Leftrightarrow f \in b}{{\wedge}E}\\
    \byrule[1]{a \in f \Leftrightarrow b \in f}{{\wedge}E}\\
    \byrule[1]{b = c}{{\wedge}E}\\
    \bydefinition[1]{\forall e: ((e \in b \Leftrightarrow e \in c) \wedge (b \in e \Leftrightarrow c \in e))}\\
    \byrule[1]{(f \in b \Leftrightarrow f \in c) \wedge (b \in f \Leftrightarrow c \in f)}{{\forall}E}\\
    \byrule[1]{f \in b \Leftrightarrow f \in c}{{\wedge}E}\\
    \byrule[1]{b \in f \Leftrightarrow c \in f}{{\wedge}E}\\
    \byrule[1]{f \in a \Leftrightarrow f \in c}{{\Leftrightarrow}T}\\
    \byrule[1]{a \in f \Leftrightarrow c \in f}{{\Leftrightarrow}T}\\
    \byrule[1]{(f \in a \Leftrightarrow f \in c) \wedge (a \in f \Leftrightarrow c \in f)}{{\wedge}I}\\
    \byrule[1]{\forall f: ((f \in a \Leftrightarrow f \in c) \wedge (a \in f \Leftrightarrow c \in f))}{{\forall}I}\\
    \bydefinition[1]{a = c}\\
    \byrule{(a = b \wedge b = c) \Rightarrow a = c}{{\Rightarrow}I}\\
    \byrule{\forall a, b, c: ((a = b \wedge b = c) \Rightarrow a = c)}{{\forall}I}
\end{align*}
\[\vdash \forall a, b, c: ((a = b \wedge b = c) \Rightarrow a = c) \quad (\text{transitivity})\]

\subsection{Substitution}
\commentary{
    By structural induction on formulas. Within each inductive case, $({=}S)$ denotes substitution already established for the immediate subformulas.
}

\begin{align*}
    \premise{a = b}\\
    \premise{a \in c}\\
    \bydefinition{\forall d: ((d \in a \Leftrightarrow d \in b) \wedge (a \in d \Leftrightarrow b \in d))}\\
    \byrule{(c \in a \Leftrightarrow c \in b) \wedge (a \in c \Leftrightarrow b \in c)}{{\forall}E}\\
    \byrule{a \in c \Leftrightarrow b \in c}{{\wedge}E}\\
    \byrule{b \in c}{{\Leftrightarrow}E}
\end{align*}
\begin{align*}
    \premise{a = b}\\
    \premise{c \in a}\\
    \bydefinition{\forall d: ((d \in a \Leftrightarrow d \in b) \wedge (a \in d \Leftrightarrow b \in d))}\\
    \byrule{(c \in a \Leftrightarrow c \in b) \wedge (a \in c \Leftrightarrow b \in c)}{{\forall}E}\\
    \byrule{c \in a \Leftrightarrow c \in b}{{\wedge}E}\\
    \byrule{c \in b}{{\Leftrightarrow}E}
\end{align*}
\[\inferrule{a = b \\ a \in c}{b \in c} \quad \inferrule{a = b \\ c \in a}{c \in b} \quad ({=}S)\]

\begin{align*}
    \premise{a = b}\\
    \premise{\forall c: P\bound{a}}\\
    \assumption[1]{P\bound{a}}\\
    \byrule[1]{P\bound{b}}{{=}S}\\
    \byrule{P\bound{a} \Rightarrow P\bound{b}}{{\Rightarrow}I}\\
    \assumption[1]{P\bound{b}}\\
    \byrule[1]{P\bound{a}}{{=}S}\\
    \byrule{P\bound{b} \Rightarrow P\bound{a}}{{\Rightarrow}I}\\
    \byrule{P\bound{a} \Leftrightarrow P\bound{b}}{{\Leftrightarrow}I}\\
    \byrule{\forall c: (P\bound{a} \Leftrightarrow P\bound{b})}{{\forall}I}\\
    \byrule{(\forall c: P\bound{a}) \Leftrightarrow (\forall c: P\bound{b})}{{\forall}D}\\
    \byrule{\forall c: P\bound{b}}{{\Leftrightarrow}E}
\end{align*}
\[\inferrule{a = b \\ \forall c: P\bound{a}}{\forall c: P\bound{b}} \quad ({=}S)\]

\begin{align*}
    \premise{a = b}\\
    \premise{P\bound{a} \nor Q}\\
    \assumption[1]{P\bound{a}}\\
    \byrule[1]{P\bound{b}}{{=}S}\\
    \byrule{P\bound{a} \Rightarrow P\bound{b}}{{\Rightarrow}I}\\
    \assumption[1]{P\bound{b}}\\
    \byrule[1]{P\bound{a}}{{=}S}\\
    \byrule{P\bound{b} \Rightarrow P\bound{a}}{{\Rightarrow}I}\\
    \byrule{P\bound{a} \Leftrightarrow P\bound{b}}{{\Leftrightarrow}I}\\
    \byrule{Q \Leftrightarrow Q}{{\Leftrightarrow}R}\\
    \byrule{(P\bound{a} \nor Q) \Leftrightarrow (P\bound{b} \nor Q)}{{\nor}D}\\
    \byrule{P\bound{b} \nor Q}{{\Leftrightarrow}E}
\end{align*}
\[\inferrule{a = b \\ P\bound{a} \nor Q}{P\bound{b} \nor Q} \quad ({=}S)\]

\[\boxed{\inferrule{a = b \\ P\bound{a}}{P\bound{b}} \quad ({=}E)}\]
